Préface

(TR)

Bu bitirme projesi, doğal dil işleme teknikleri kullanılarak geliştirilen çok dilli metin özetleme modellerinin farklı diller üzerindeki performansını incelemek amacıyla hazırlanmıştır. Çalışma kapsamında, mBART tabanlı modeller XLSum veri kümesi kullanılarak 15 farklı dilde haber metinlerinin özetlenmesi için eğitilmiş ve değerlendirilmiştir. Proje, Galatasaray Üniversitesi Bilgisayar Mühendisliği lisans programı kapsamında gerçekleştirilmiş olup, teorik bilgilerin gerçek dünya problemlerine uygulanmasına olanak sağlamıştır.

Bu konunun seçilmesindeki temel motivasyon, lisans eğitimim süresince büyük dil modellerinin hızla gelişmesi ve son kullanıcı uygulamalarında yaygınlaşmasıyla birlikte, insan diliyle olan etkileşimlerinin giderek daha önemli hale gelmesidir. Proje sürecinde veri ön işleme, model eğitimi, çok dilli değerlendirme ve sonuçların analizi gibi aşamalarda önemli deneyimler elde ettim.

""" Ayrıca, geliştirilen modeller bir web uygulaması ile entegre edilerek RSS kaynaklarından alınan güncel haberlerin otomatik olarak özetlenmesi sağlanmıştır. Bu uygulama sayesinde kullanıcılar, farklı kaynaklardan gelen haber özetlerine ve ilgili bağlantılara akış halinde erişebilmektedir. """

Bu çalışmanın gerçekleştirilmesinde değerli rehberliği ve destekleri için danışmanım Dr. Öğr. Üyesi İsmail Burak Parlak’a teşekkür ederim. Ayrıca, süreç boyunca desteklerini esirgemeyen aileme ve arkadaşlarıma da teşekkürlerimi sunarım.

Bu çalışmanın, konuya ilgi duyan okuyucular için faydalı olmasını temenni ederim.

---

(FR)

Ce projet de fin d’études a été réalisé afin d’étudier les performances des modèles de résumé de texte multilingues développés à l’aide des techniques de traitement automatique du langage naturel. Dans ce cadre, des modèles basés sur mBART ont été entraînés et évalués pour résumer des articles de presse dans 15 langues différentes en utilisant le jeu de données XLSum. Ce travail a été effectué dans le cadre du programme de licence en ingénierie informatique de l’Université Galatasaray et a permis d’appliquer des connaissances théoriques à un problème concret.

Le choix de ce sujet est motivé par le développement rapide des grands modèles de langage et par leur utilisation de plus en plus fréquente dans les applications destinées aux utilisateurs. Cette évolution montre une interaction de plus en plus importante avec le langage humain. Au cours de ce projet, j’ai acquis des compétences importantes en traitement des données, en entraînement de modèles, en évaluation multilingue et en analyse des résultats.

""" De plus, les modèles développés ont été intégrés dans une application web qui permet de résumer automatiquement des actualités provenant de flux RSS. Grâce à cette application, les utilisateurs peuvent accéder à un flux continu de résumés d’articles ainsi qu’aux liens vers les sources originales. """

Je tiens à remercier mon encadrant, Dr. İsmail Burak Parlak, pour son accompagnement et ses précieux conseils. Je remercie également ma famille et mes amis pour leur soutien tout au long de ce projet.

J’espère que ce travail sera utile pour les lecteurs intéressés par ce domaine.